\chapter{Les acides carboxyliques et leurs dérivés}

\textit{Les acides carboxyliques constituent une famille de composés organiques très répandus dans la nature (acide acétique du vinaigre, acide citrique des agrumes, acides gras des lipides). Leur groupe fonctionnel, le groupe carboxyle \ce{-COOH}, leur confère une acidité caractéristique et une grande réactivité. Ils sont à l'origine de nombreuses transformations chimiques permettant d'obtenir des dérivés fonctionnels tels que les esters, les anhydrides d'acide, les chlorures d'acyle et les amides. Ce chapitre explore la structure, la nomenclature, les propriétés et les réactions de ces composés essentiels.}

\section{Structure et nomenclature des acides carboxyliques}

\begin{definition}[Groupe carboxyle]
Un acide carboxylique est un composé organique qui possède le groupe fonctionnel carboxyle \ce{-COOH}, résultant de l'association d'un groupe carbonyle \ce{C=O} et d'un groupe hydroxyle \ce{-OH} portés par le même atome de carbone. Sa formule générale est \ce{R-COOH} (ou \ce{Ar-COOH} pour les acides aromatiques).
\end{definition}

\begin{illustration}[Formule développée de l'acide carboxylique]
\begin{center}
\chemfig{R-C(=[1]O)-[7]OH}
\end{center}
\fig{Groupe carboxyle : association d'un carbonyle et d'un hydroxyle.}
\end{illustration}

\subsection{Nomenclature systématique (IUPAC)}

\begin{propriete}[Règles de nomenclature]
\begin{itemize}
    \item On repère la chaîne carbonée la plus longue contenant le groupe \ce{-COOH}.
    \item Le nom de l'acide dérive de l'alcane correspondant en remplaçant le \textit{-e} final par \textit{-oïque}, précédé du mot \textit{acide}.
    \item Le carbone du groupe \ce{-COOH} est toujours le numéro 1.
    \item Les substituants sont nommés et numérotés comme d'habitude.
\end{itemize}
\end{propriete}

\begin{exemple}
\begin{itemize}
    \item \ce{HCOOH} : acide méthanoïque (acide formique)
    \item \ce{CH3COOH} : acide éthanoïque (acide acétique)
    \item \ce{CH3CH2CH2COOH} : acide butanoïque
    \item \ce{CH3CH(Cl)COOH} : acide 2-chloropropanoïque
    \item \chemfig{*6(-=-(-COOH)=-=)} : acide benzoïque
\end{itemize}
\end{exemple}

\subsection{Nomenclature courante}

\begin{remarque}
De nombreux acides carboxyliques possèdent des noms usuels historiques, souvent dérivés de leur source naturelle :
\begin{itemize}
    \item Acide formique (fourmis, latin \textit{formica})
    \item Acide acétique (vinaigre, latin \textit{acetum})
    \item Acide butyrique (beurre rance, latin \textit{butyrum})
    \item Acide palmitique (huile de palme)
    \item Acide stéarique (suif, grec \textit{stear})
\end{itemize}
\end{remarque}

\section{Propriétés physiques des acides carboxyliques}

\begin{propriete}[Propriétés physiques]
\begin{itemize}
    \item Les acides carboxyliques sont polaires et forment des liaisons hydrogène intermoléculaires (dimères).
    \item Leurs points d'ébullition sont élevés, supérieurs à ceux des alcools de masse molaire comparable.
    \item Les premiers termes (\ce{C1} à \ce{C4}) sont miscibles à l'eau ; la solubilité diminue quand la chaîne carbonée s'allonge.
    \item Les acides carboxyliques ont une odeur souvent piquante (acide acétique) ou désagréable (acide butyrique).
\end{itemize}
\end{propriete}

\begin{illustration}[Dimère d'acide carboxylique]
\begin{center}
\chemfig{R-C(=[1]O)-[7]OH...O=[6]C(-[2]R)-[7]OH}
\end{center}
\fig{Formation de dimères par liaisons hydrogène.}
\end{illustration}

\section{Acidité des acides carboxyliques}

\subsection{Équilibre acido-basique}

\begin{loi}[Constante d'acidité]
Un acide carboxylique \ce{RCOOH} se comporte comme un acide faible en solution aqueuse :
\[
\ce{RCOOH + H2O <=> RCOO- + H3O+}
\]
La constante d'acidité est définie par :
\[
K_a = \frac{[\ce{RCOO-}][\ce{H3O+}]}{[\ce{RCOOH}]}
\]
et le \textit{pKa} par : $\text{p}K_a = -\log K_a$.
\end{loi}

\begin{propriete}[Valeurs de pKa]
Les acides carboxyliques ont des pKa compris entre 4 et 5 (plus acides que l'eau, moins acides que les acides minéraux forts).
\begin{itemize}
    \item Acide éthanoïque : pKa = 4,76
    \item Acide méthanoïque : pKa = 3,75
    \item Acide benzoïque : pKa = 4,20
    \item Acide chloroéthanoïque : pKa = 2,86
\end{itemize}
\end{propriete}

\begin{remarque}
L'acidité est due à la stabilisation par résonance de la base conjuguée (ion carboxylate \ce{RCOO-}), dont la charge négative est délocalisée sur les deux atomes d'oxygène.
\end{remarque}

\begin{illustration}[Résonance de l'ion carboxylate]
\begin{center}
\chemfig{R-C(=[1]O)-[7]O-} \hspace{1cm} \chemfig{R-C(-[1]O-)-[7]O}
\end{center}
\fig{Délocalisation de la charge négative dans l'ion carboxylate.}
\end{illustration}

\subsection{Facteurs influençant l'acidité}

\begin{propriete}[Effets électroniques]
\begin{itemize}
    \item Les groupements attracteurs d'électrons (halogènes, \ce{-NO2}, \ce{-CN}) augmentent l'acidité (pKa plus faible).
    \item Les groupements donneurs d'électrons (alkyles) diminuent l'acidité (pKa plus élevé).
    \item L'effet est d'autant plus fort que le substituant est proche du groupe carboxyle.
\end{itemize}
\end{propriete}

\begin{exemple}
Comparaison des pKa :
\begin{itemize}
    \item \ce{CH3COOH} : pKa = 4,76
    \item \ce{ClCH2COOH} : pKa = 2,86
    \item \ce{Cl2CHCOOH} : pKa = 1,29
    \item \ce{Cl3CCOOH} : pKa = 0,66
\end{itemize}
\end{exemple}

\section{Obtention des acides carboxyliques}

\subsection{Oxydation des alcools primaires}

\begin{experience}[Oxydation ménagée]
Un alcool primaire peut être oxydé en acide carboxylique par un oxydant fort comme le permanganate de potassium \ce{KMnO4} en milieu acide ou le dichromate de potassium \ce{K2Cr2O7}.
\[
\ce{RCH2OH + 2[O] -> RCOOH + H2O}
\]
\end{experience}

\begin{exemple}
Oxydation de l'éthanol en acide éthanoïque :
\[
\ce{CH3CH2OH + 2[O] ->[K2Cr2O7/H2SO4] CH3COOH + H2O}
\]
\end{exemple}

\subsection{Oxydation des aldéhydes}

\begin{propriete}
Les aldéhydes sont facilement oxydés en acides carboxyliques par des oxydants doux (réactif de Tollens, liqueur de Fehling) ou forts.
\[
\ce{RCHO + [O] -> RCOOH}
\]
\end{propriete}

\subsection{Hydrolyse des dérivés d'acide}

\begin{remarque}
Les esters, les chlorures d'acyle, les anhydrides d'acide et les amides peuvent être hydrolysés pour donner l'acide carboxylique correspondant (voir section \ref{sec:derives}).
\end{remarque}

\section{Réactions des acides carboxyliques}

\subsection{Réaction avec les bases : formation de sels}

\begin{loi}[Neutralisation]
Un acide carboxylique réagit avec une base forte (soude, potasse) pour former un sel carboxylate et de l'eau :
\[
\ce{RCOOH + NaOH -> RCOONa + H2O}
\]
\end{loi}

\begin{exemple}
Réaction de l'acide éthanoïque avec la soude :
\[
\ce{CH3COOH + NaOH -> CH3COONa + H2O}
\]
Le sel obtenu est l'éthanoate de sodium (acétate de sodium).
\end{exemple}

\begin{propriete}[Propriétés des sels carboxylates]
\begin{itemize}
    \item Les carboxylates de sodium ou de potassium sont solubles dans l'eau.
    \item Ce sont des composés ioniques, contrairement aux acides carboxyliques.
    \item Ils sont utilisés comme savons (carboxylates à longue chaîne).
\end{itemize}
\end{propriete}

\subsection{Réaction avec les alcools : estérification}

\begin{definition}[Estérification]
L'estérification est la réaction entre un acide carboxylique et un alcool, en milieu acide, pour former un ester et de l'eau.
\[
\ce{RCOOH + R$'$OH <=>[H+] RCOOR$'$ + H2O}
\]
\end{definition}

\begin{experience}[Synthèse de l'éthanoate d'éthyle]
\begin{itemize}
    \item Dans un ballon, introduire \SI{10}{\mL} d'acide éthanoïque, \SI{10}{\mL} d'éthanol et quelques gouttes d'acide sulfurique concentré.
    \item Chauffer à reflux pendant environ 30 minutes.
    \item L'ester formé (éthanoate d'éthyle) est reconnaissable à son odeur fruitée caractéristique (poire).
\end{itemize}
\end{experience}

\begin{propriete}[Caractéristiques de l'estérification]
\begin{itemize}
    \item Réaction lente et limitée (équilibre chimique).
    \item Catalysée par les ions \ce{H+} (acide sulfurique concentré).
    \item Réaction athermique (faible variation d'enthalpie).
    \item Le rendement peut être amélioré en éliminant l'eau ou en utilisant un excès d'un réactif.
\end{itemize}
\end{propriete}

\begin{illustration}[Mécanisme de l'estérification]
\begin{center}
\chemfig{R-C(=[1]O)-[7]OH} + \chemfig{R'-OH} \hspace{0.5cm} \chemfig{R-C(=[1]O)-[7]O-R'} + \chemfig{H_2O}
\end{center}
\fig{Équation bilan de l'estérification.}
\end{illustration}

\begin{loi}[Constante d'équilibre]
Pour l'estérification :
\[
K = \frac{[\text{ester}][\ce{H2O}]}{[\text{acide}][\text{alcool}]}
\]
À température ambiante, $K \approx 4$ pour un acide et un alcool primaire.
\end{loi}

\begin{exemple}
Calcul du rendement : si on mélange \SI{1}{\mol} d'acide et \SI{1}{\mol} d'alcool, à l'équilibre on obtient environ $\frac{2}{3}$ de mole d'ester (rendement $\approx 67\%$).
\end{exemple}

\subsection{Réduction des acides carboxyliques}

\begin{propriete}
Les acides carboxyliques sont difficiles à réduire. On utilise des hydrures puissants comme \ce{LiAlH4} (aluminohydrure de lithium) pour les transformer en alcools primaires :
\[
\ce{RCOOH + 4[H] ->[LiAlH4] RCH2OH + H2O}
\]
\end{propriete}

\section{Dérivés d'acide carboxylique}
\label{sec:derives}

\begin{definition}[Dérivés fonctionnels]
Les dérivés d'acide carboxylique sont des composés obtenus en remplaçant le groupe \ce{-OH} du carboxyle par un autre groupe (halogène, \ce{-OR$'$}, \ce{-OCOR$'$}, \ce{-NH2}, etc.).
\end{definition}

\subsection{Les chlorures d'acyle}

\begin{propriete}[Chlorures d'acyle]
Formule générale : \ce{RCOCl}. Obtenus par action du chlorure de thionyle \ce{SOCl2} ou du pentachlorure de phosphore \ce{PCl5} sur l'acide carboxylique :
\[
\ce{RCOOH + SOCl2 -> RCOCl + SO2 + HCl}
\]
\end{propriete}

\begin{remarque}
Les chlorures d'acyle sont très réactifs et utilisés comme agents d'acylation.
\end{remarque}

\subsection{Les anhydrides d'acide}

\begin{propriete}[Anhydrides d'acide]
Formule générale : \ce{RCO-O-COR$'$}. Obtenus par déshydratation de deux acides carboxyliques (souvent avec \ce{P2O5}) :
\[
\ce{RCOOH + R$'$COOH ->[P2O5] RCO-O-COR$'$ + H2O}
\]
\end{propriete}

\begin{exemple}
Anhydride éthanoïque : \ce{(CH3CO)2O}, obtenu à partir de deux molécules d'acide éthanoïque.
\end{exemple}

\subsection{Les esters}

\begin{definition}[Ester]
Un ester est un composé de formule générale \ce{RCOOR$'$}, où \ce{R} provient de l'acide et \ce{R$'$} de l'alcool.
\end{definition}

\begin{illustration}[Formule générale d'un ester]
\begin{center}
\chemfig{R-C(=[1]O)-[7]O-R'}
\end{center}
\fig{Groupe fonctionnel ester.}
\end{illustration}

\subsubsection{Nomenclature des esters}

\begin{propriete}[Règles de nomenclature]
Le nom d'un ester se construit à partir du nom de l'acide (en remplaçant \textit{-oïque} par \textit{-oate}) suivi du nom du groupe alkyle provenant de l'alcool.
\end{propriete}

\begin{exemple}
\begin{itemize}
    \item \ce{CH3COOCH2CH3} : éthanoate d'éthyle (acétate d'éthyle)
    \item \ce{HCOOCH3} : méthanoate de méthyle (formiate de méthyle)
    \item \ce{C6H5COOCH2CH3} : benzoate d'éthyle
\end{itemize}
\end{exemple}

\subsubsection{Réactions des esters}

\begin{propriete}[Hydrolyse acide]
L'hydrolyse acide d'un ester est la réaction inverse de l'estérification :
\[
\ce{RCOOR$'$ + H2O <=>[H+] RCOOH + R$'$OH}
\]
C'est une réaction lente et limitée.
\end{propriete}

\begin{propriete}[Saponification]
La saponification est l'hydrolyse basique d'un ester par une base forte (soude ou potasse) :
\[
\ce{RCOOR$'$ + NaOH -> RCOONa + R$'$OH}
\]
\end{propriete}

\begin{experience}[Saponification]
\begin{itemize}
    \item Dans un ballon, introduire \SI{5}{\mL} d'éthanoate d'éthyle et \SI{20}{\mL} de soude \SI{2}{\mol\per\litre}.
    \item Chauffer à reflux pendant 20 minutes.
    \item L'ester disparaît ; on obtient une solution d'éthanoate de sodium et d'éthanol.
    \item La réaction est totale et irréversible (contrairement à l'hydrolyse acide).
\end{itemize}
\end{experience}

\begin{remarque}
La saponification est utilisée industriellement pour fabriquer des savons à partir de corps gras (esters d'acides gras et de glycérol).
\end{remarque}

\subsection{Les amides}

\begin{propriete}[Amides]
Formule générale : \ce{RCONH2} (amide primaire), \ce{RCONHR$'$} (amide secondaire), \ce{RCONR$'$R$'$$'$} (amide tertiaire).
Obtenus par action de l'ammoniac ou d'une amine sur un chlorure d'acyle ou un anhydride :
\[
\ce{RCOCl + 2NH3 -> RCONH2 + NH4Cl}
\]
\end{propriete}

\begin{exemple}
\ce{CH3COCl + 2NH3 -> CH3CONH2 + NH4Cl} (éthanamide)
\end{exemple}

\section{Synthèse d'un ester par différentes voies}

\begin{propriete}[Voies de synthèse d'un ester]
Un ester \ce{RCOOR$'$} peut être obtenu par plusieurs méthodes :
\begin{enumerate}
    \item \textbf{Estérification directe} : \ce{RCOOH + R$'$OH <=>[H+] RCOOR$'$ + H2O}
    \item \textbf{Action d'un chlorure d'acyle} : \ce{RCOCl + R$'$OH -> RCOOR$'$ + HCl}
    \item \textbf{Action d'un anhydride d'acide} : \ce{(RCO)2O + R$'$OH -> RCOOR$'$ + RCOOH}
    \item \textbf{Action d'un sel carboxylate} : \ce{RCOONa + R$'$Cl -> RCOOR$'$ + NaCl}
\end{enumerate}
\end{propriete}

\begin{exemple}
Synthèse de l'éthanoate d'éthyle par trois méthodes :
\begin{itemize}
    \item \ce{CH3COOH + CH3CH2OH <=>[H+] CH3COOCH2CH3 + H2O}
    \item \ce{CH3COCl + CH3CH2OH -> CH3COOCH2CH3 + HCl}
    \item \ce{(CH3CO)2O + CH3CH2OH -> CH3COOCH2CH3 + CH3COOH}
\end{itemize}
\end{exemple}

\begin{remarque}
Les méthodes utilisant les chlorures d'acyle ou les anhydrides sont plus rapides et donnent de meilleurs rendements que l'estérification directe, mais nécessitent des réactifs plus coûteux et moins accessibles.
\end{remarque}

\section{L'essentiel à retenir}

\begin{synthese}
\subsection*{Acides carboxyliques}
\begin{itemize}
    \item Groupe fonctionnel : \ce{-COOH} (carboxyle)
    \item Nomenclature : acide \textit{-oïque}
    \item Acidité : pKa $\approx$ 4-5, base conjuguée stabilisée par résonance
    \item Obtention : oxydation d'alcools primaires ou d'aldéhydes
    \item Réactions : neutralisation (sels), estérification (esters), réduction (alcools)
\end{itemize}

\subsection*{Estérification}
\begin{itemize}
    \item \ce{RCOOH + R$'$OH <=>[H+] RCOOR$'$ + H2O}
    \item Réaction lente, limitée, catalysée par \ce{H+}
    \item Constante d'équilibre $K \approx 4$ (à température ambiante)
    \item Rendement amélioré par excès d'un réactif ou élimination d'eau
\end{itemize}

\subsection*{Dérivés d'acide}
\begin{itemize}
    \item Chlorures d'acyle : \ce{RCOCl} (très réactifs)
    \item Anhydrides : \ce{(RCO)2O} (réactifs)
    \item Esters : \ce{RCOOR$'$} (hydrolyse acide limitée, saponification totale)
    \item Amides : \ce{RCONH2} (moins réactifs)
\end{itemize}

\subsection*{Saponification}
\begin{itemize}
    \item \ce{RCOOR$'$ + NaOH -> RCOONa + R$'$OH}
    \item Réaction totale et irréversible
    \item Utilisée pour la fabrication des savons
\end{itemize}

\subsection*{Pièges à éviter}
\begin{itemize}
    \item Ne pas confondre estérification (acide + alcool) et saponification (ester + base)
    \item L'hydrolyse acide d'un ester est l'inverse de l'estérification (équilibre)
    \item La saponification est totale car le carboxylate formé ne réagit pas avec l'alcool
    \item Les chlorures d'acyle sont plus réactifs que les anhydrides, eux-mêmes plus réactifs que les acides
\end{itemize}
\end{synthese}

\section{Méthodes \& savoir-faire}

\methode{Reconnaître et nommer un acide carboxylique}
Pour nommer un acide carboxylique :
\begin{enumerate}
    \item Identifier la chaîne carbonée la plus longue contenant le groupe \ce{-COOH}.
    \item Numéroter à partir du carbone du groupe \ce{-COOH} (carbone 1).
    \item Remplacer le \textit{-e} final de l'alcane par \textit{-oïque}, précédé de \textit{acide}.
    \item Indiquer la position des substituants.
\end{enumerate}

\begin{exemple}
Nommer \ce{CH3CH(Cl)CH2COOH}.
\begin{itemize}
    \item Chaîne principale : 4 carbones (butane)
    \item Groupe \ce{-COOH} en position 1
    \item Chlore en position 3
    \item Nom : acide 3-chlorobutanoïque
\end{itemize}
\end{exemple}

\methode{Écrire l'équation d'estérification}
Pour écrire l'équation d'estérification entre un acide carboxylique et un alcool :
\begin{enumerate}
    \item Identifier l'acide \ce{RCOOH} et l'alcool \ce{R$'$OH}.
    \item L'ester formé est \ce{RCOOR$'$}.
    \item L'autre produit est l'eau \ce{H2O}.
    \item Équilibrer l'équation (déjà équilibrée : 1 acide + 1 alcool $\rightarrow$ 1 ester + 1 eau).
\end{enumerate}

\begin{exemple}
Écrire l'équation d'estérification entre l'acide propanoïque et le méthanol.
\begin{itemize}
    \item Acide propanoïque : \ce{CH3CH2COOH}
    \item Méthanol : \ce{CH3OH}
    \item Ester : \ce{CH3CH2COOCH3} (propanoate de méthyle)
    \item Équation : \ce{CH3CH2COOH + CH3OH <=>[H+] CH3CH2COOCH3 + H2O}
\end{itemize}
\end{exemple}

\methode{Calculer le rendement d'une estérification}
Pour calculer le rendement d'une estérification à partir des quantités initiales et de la constante d'équilibre :
\begin{enumerate}
    \item Écrire l'équation et le tableau d'avancement.
    \item Exprimer la constante d'équilibre $K$ en fonction de l'avancement $x$.
    \item Résoudre l'équation pour trouver $x$.
    \item Le rendement est $r = \frac{x}{n_{\text{initial}}}$ (pour le réactif limitant).
\end{enumerate}

\begin{exemple}
On mélange \SI{2.0}{\mol} d'acide éthanoïque et \SI{1.0}{\mol} d'éthanol. $K = 4$. Calculer le rendement.
\begin{itemize}
    \item Tableau : initial (2, 1, 0, 0), équilibre ($2-x$, $1-x$, $x$, $x$)
    \item $K = \frac{x^2}{(2-x)(1-x)} = 4$
    \item $x^2 = 4(2-x)(1-x) = 4(2 - 3x + x^2) = 8 - 12x + 4x^2$
    \item $0 = 8 - 12x + 3x^2$ soit $3x^2 - 12x + 8 = 0$
    \item $\Delta = 144 - 96 = 48$, $\sqrt{\Delta} \approx 6.93$
    \item $x = \frac{12 - 6.93}{6} \approx 0.845$ (l'autre racine $>1$)
    \item Rendement par rapport à l'éthanol (limitant) : $r = \frac{0.845}{1.0} = 84.5\%$
\end{itemize}
\end{exemple}

\methode{Distinguer hydrolyse acide et saponification}
Pour distinguer ces deux réactions :
\begin{enumerate}
    \item L'hydrolyse acide utilise un acide fort (\ce{H+}) et donne l'acide carboxylique et l'alcool (réaction limitée).
    \item La saponification utilise une base forte (\ce{OH-}) et donne le sel carboxylate et l'alcool (réaction totale).
\end{enumerate}

\begin{exemple}
Pour l'ester \ce{CH3COOCH2CH3} :
\begin{itemize}
    \item Hydrolyse acide : \ce{CH3COOCH2CH3 + H2O <=>[H+] CH3COOH + CH3CH2OH}
    \item Saponification : \ce{CH3COOCH2CH3 + NaOH -> CH3COONa + CH3CH2OH}
\end{itemize}
\end{exemple}

\methode{Comparer la réactivité des dérivés d'acide}
Pour classer les dérivés d'acide par réactivité décroissante :
\begin{enumerate}
    \item Chlorures d'acyle (plus réactifs)
    \item Anhydrides d'acide
    \item Esters
    \item Amides (moins réactifs)
\end{enumerate}
Cette réactivité est liée à la force du groupe partant.

\begin{exemple}
Pour obtenir un ester \ce{RCOOR$'$} à partir de \ce{RCOOH} :
\begin{itemize}
    \item Avec \ce{RCOCl} : réaction rapide et totale à température ambiante
    \item Avec \ce{(RCO)2O} : réaction rapide, nécessite parfois un catalyseur
    \item Avec \ce{RCOOH} directement : réaction lente et limitée, nécessite chauffage et catalyseur
\end{itemize}
\end{exemple}

\methode{Identifier un ester par son odeur}
Pour identifier un ester en laboratoire :
\begin{enumerate}
    \item Les esters ont des odeurs fruitées caractéristiques (poire, banane, ananas, pomme).
    \item L'éthanoate d'éthyle sent la poire.
    \item L'éthanoate d'isoamyle sent la banane.
    \item Le butanoate d'éthyle sent l'ananas.
\end{enumerate}

\begin{exemple}
Un liquide incolore à l'odeur de banane est probablement l'éthanoate d'isoamyle \ce{CH3COOCH2CH2CH(CH3)2}.
\end{exemple}

\methode{Écrire le mécanisme d'estérification}
Pour écrire le mécanisme d'estérification en milieu acide :
\begin{enumerate}
    \item Protonation du carbonyle de l'acide.
    \item Attaque nucléophile de l'alcool sur le carbone carbonylé.
    \item Transfert de proton.
    \item Départ d'eau.
    \item Déprotonation pour donner l'ester.
\end{enumerate}

\begin{exemple}
Mécanisme simplifié de l'estérification de l'acide éthanoïque par l'éthanol :
\begin{itemize}
    \item \ce{CH3COOH + H+ <=> CH3C(OH)2+}
    \item \ce{CH3C(OH)2+ + CH3CH2OH -> CH3C(OH)(OCH2CH3)+ + H2O}
    \item \ce{CH3C(OH)(OCH2CH3)+ -> CH3COOCH2CH3 + H+}
\end{itemize}
\end{exemple}

\methode{Calculer le pH d'une solution d'acide carboxylique}
Pour calculer le pH d'une solution d'acide carboxylique de concentration $C$ :
\begin{enumerate}
    \item Vérifier que $C \gg K_a$ (approximation valide si $C > 100 K_a$).
    \item Utiliser la formule : $[\ce{H3O+}] = \sqrt{K_a \times C}$.
    \item pH $= -\log[\ce{H3O+}]$.
    \item Si l'approximation n'est pas valide, résoudre l'équation du second degré.
\end{enumerate}

\begin{exemple}
Calculer le pH d'une solution d'acide éthanoïque à \SI{0.10}{\mol\per\litre} (pKa = 4,76).
\begin{itemize}
    \item $K_a = 10^{-4.76} \approx 1.74 \times 10^{-5}$
    \item $C = 0.10 > 100 \times 1.74 \times 10^{-5} = 1.74 \times 10^{-3}$ (approximation valide)
    \item $[\ce{H3O+}] = \sqrt{1.74 \times 10^{-5} \times 0.10} = \sqrt{1.74 \times 10^{-6}} \approx 1.32 \times 10^{-3}$
    \item pH $= -\log(1.32 \times 10^{-3}) \approx 2.88$
\end{itemize}
\end{exemple}